\DocumentMetadata{
  lang = pt-BR,
  pdfstandard = A-2b,
  pdfversion = 1.7
}

\documentclass{abntexto-ufc}

\ufcsetup{
  type = article,
  print-mode = double-sided,
  cover = true,
  catalog-card = false,
  font = times,
  strict-font = false,
  author = {UFCARTICLEAUTHORMARKER\footnote{Pesquisador do Centro de Teste da Universidade Federal do Ceará. Contato institucional: autor@example.org.}},
  title = {UFCARTICLETITLEMARKER},
  subtitle = {Scientific article runtime validation},
  title-variant = {UFCARTICLEOTHERTITLEMARKER},
  location = {Fortaleza},
  year = {2026},
  approval-date = {20 de agosto de 2026}
}

\date{18 de agosto de 2026}
\ufcAddBibliographyResource{tests/fixtures/referencias-v2.bib}

\begin{document}
\pretextual

\ufcPrintCover
\ufcPrintTitlePage
\ufcPrintApprovalPage
\ufcPrintSummary{tests/fixtures/artigo-resumo}
\ufcPrintAbstract{tests/fixtures/artigo-abstract}
\ufcPrintTableOfContents

\textual
\section{Introdução}
UFCARTICLEINTROMARKER apresenta o contexto do problema, o objetivo do trabalho e a organização da análise. Este parágrafo confirma o recuo de primeira linha e o fluxo textual em espaço simples.

\newpage
UFCARTICLEFLOWMARKERBEFOREDEVELOPMENT.
\section{Desenvolvimento}
UFCARTICLEDEVELOPMENTMARKER descreve os fundamentos, o método e os resultados necessários à validação do perfil científico. O conteúdo permanece no mesmo fluxo da seção primária anterior, sem quebra automática introduzida pelo comando de seção.

UFCARTICLEFLOWMARKERBEFOREFINAL.
\section{Considerações finais}
UFCARTICLEFINALMARKER sintetiza os resultados observados e registra as conclusões do artigo. A seção deve permanecer na mesma página do marcador imediatamente anterior quando houver espaço físico suficiente.

\nocite{silva2020}
\ufcPrintReferences

\end{document}
